\input{../tex-inputs/latex-leseansicht-vorspann}

\section[Arthur Schnitzler an Marie Holzer, 25. 4. 1908]{L04329 Arthur Schnitzler an Marie Holzer, 25. 4. 1908}
\nopagebreak\mylabel{L04329v}
\rehead{ }\normalsize\beginnumbering\briefempfaengerindex{Holzer, Marie@\textsc{Holzer, Marie}!zzzSchnitzler, Arthur@\emph{von Arthur Schnitzler}!1908-04-251@{25. 4. 1908}|(be}
\toendnotes[C]{\smallbreak\pagebreak[2]}
\correspDesc{Versand  durch Arthur Schnitzler am 25. 4. 1908 in Wien
\newline{}Erhalt  durch Marie Holzer im Zeitraum [26. 4. 1908 – 30. 4. 1908?] in Prag}\toendnotes[C]{\smallbreak}
\Standort{DLA, A:Schnitzler, HS.1985.1.1042.}
\physDesc{Brief, Durchschlag, 1 Blatt, 1 Seite, 454 Zeichen
\newline{}Schreibmaschine
\newline{}Handschrift: 1) roter Buntstift (\noindent{}zwei Unterstreichungen)\hspace{1em}2) Bleistift, lateinische Kurrent (\noindent{}Vermerk »Holzer«)\hspace{1em}}\toendnotes[C]{\smallbreak}
\pstart
           \raggedleft{}{\pb}25. April 08.\pend
           
\pstart{}Verehrte gnädige Frau,\pend\vspace{0.5em}
\pstart
           Ich danke Ihnen herzlich für Ihren \label{K_L04221-1v}\edtext{Brief}{\lemma{\textnormal{\emph{Brief}}}\Cendnote{\textnormal{XXXX Auszeichnungsfehler: Dokument L04220 nicht gefunden.}}}\label{K_L04221-1} und für das Interesse, das Sie an der kleinen Novelle\pwindex{Schnitzler, Arthur 15. 5. 1862 Wien – 21. 10. 1931 ebd.@\textsc{Schnitzler, Arthur} (15. 5. 1862 Wien – 21. 10. 1931 ebd.), \emph{Schriftsteller, Mediziner}!Tod des Junggesellen. Novelle@\strich\emph{Der Tod des Junggesellen. Novelle}|pwv} nehmen. Dass ich Sie trotzdem
          bitten muss auf einen Kom{[}m{]}entar zu verzichten wird kaum jemand besser
          zu verstehen im Stande sein, als Sie gnädige Frau, die schon mehr als einmal eine so
          erfreuliche Einsicht das Wesen dichterischen Schaffens bewiesen hat.\pend
           
\pstart
           Mit verbindlichsten Grüssen{\\[\baselineskip]} Ihr aufrichtig ergebener\pend
           \leftskip=0em{}{\vspace{1\baselineskip}}
\pstart
           \noindent{}Frau Marie Holzer, Prag\oindex{Prag@\textbf{Prag}, \emph{Land}|pw}.\pend
           \selectlanguage{ngerman}\endnumbering\briefempfaengerindex{Holzer, Marie@\textsc{Holzer, Marie}!zzzSchnitzler, Arthur@\emph{von Arthur Schnitzler}!1908-04-251@{25. 4. 1908}|)be}\mylabel{L04329h}
\begin{anhang}
\end{anhang}\newcommand{\dateiname}{L04329}\newcommand{\titel}{Arthur Schnitzler an Marie Holzer, 25. 4. 1908}\newcommand{\editorInnen}{Herausgegeben von Jahnke, SelmaMüller, Martin Anton}\input{../tex-inputs/latex-leseansicht-abspann}
